\setcounter{chapter}{0}
\setcounter{section}{0}
\setcounter{table}{0}
\setcounter{figure}{0}
%\addcontentsline{toc}{chapter}{CƠ SỞ LÝ THUYẾT}
\chapter{CƠ SỞ LÝ THUYẾT}

\textit{Chương này trình bày những công nghệ chính sử dụng trong đồ án.}

\section{Tổng quan đề tài}
Sinh viên trình bày tổng quan về đề tài, bối cảnh ứng dụng và vấn đề cần giải quyết.

\section{Cơ sở lý thuyết}
Sinh viên trình bày các khái niệm, mô hình, thuật toán, phương pháp hoặc công nghệ nền tảng liên quan đến đề tài.
Ví dụ, khi trình bày tài liệu nền tảng, sinh viên có thể trích dẫn như sau: quy trình phát triển phần mềm được trình bày trong tài liệu của Sommerville~\cite{sommerville2016software}.

Ví dụ danh sách gạch đầu dòng:
\begin{itemize}
    \item Công nghệ hoặc nền tảng thứ nhất sử dụng trong đồ án.
    \item Công nghệ hoặc nền tảng thứ hai sử dụng trong đồ án.
    \item Công nghệ hoặc nền tảng thứ ba sử dụng trong đồ án.
\end{itemize}

\section{Các công trình hoặc hệ thống liên quan}
Sinh viên tóm tắt các nghiên cứu, sản phẩm hoặc hệ thống tương tự. Nên nêu ưu điểm, hạn chế và bài học kế thừa.

\section{Định hướng giải quyết}
Sinh viên nêu hướng tiếp cận dự kiến, các giả định ban đầu và lý do lựa chọn phương án thực hiện.


